\section{Ideas extras creativas}

\subsection{Velocidad de escape local durante el encuentro}

Calcular a lo largo del tiempo la velocidad de escape local del sistema Tierra--Apophis,
\[
v_{esc} = \sqrt{\frac{2 G M_{\oplus}}{r}},
\]
y compararla con la velocidad relativa de Apophis. Mostrar que en ningún momento $v_{rel} < v_{esc}$ y explicar por qué el asteroide no queda capturado.

\subsection{Comparar con la segunda ley de Kepler}

Calcular numéricamente el área barrida por el vector relativo Sol--Apophis en intervalos de tiempo iguales. Verificar que dichas áreas son iguales en la aproximación de 2 cuerpos y estudiar cómo esta ley se viola ligeramente cuando se incluyen perturbaciones planetarias.

\subsection{Energía de la hipérbola geocéntrica}

Calcular la energía específica geocéntrica,
\[
\mathcal{E}_{geo} = \frac{v^2}{2} - \frac{G M_{\oplus}}{r},
\]
a lo largo de la trayectoria de Apophis. Mostrar que permanece positiva, como corresponde a una trayectoria hiperbólica no capturada, y que alcanza un mínimo cerca del instante de máxima aproximación.

\subsection{Posición en el cielo desde Colombia}

Convertir las coordenadas eclípticas de Apophis en el momento del encuentro a ascensión recta y declinación. Determinar en qué constelación estará y evaluar si será visible desde Medellín durante la noche del 13 de abril de 2029.

\subsection{Intercambio de energía por marco rotante Tierra--Apophis}

Construir la dinámica en el marco rotante del sistema Tierra--Apophis durante las horas cercanas al encuentro y calcular, paso a paso, la variación de energía efectiva del asteroide. Identificar en qué tramo exacto gana energía heliocéntrica y en qué tramo la pierde. Relacionar ese resultado con un ``asistencia gravitacional natural'' y estimar cómo cambia su próximo período orbital alrededor del Sol.

\subsection{Mapa de sensibilidad de la distancia mínima}

Generar un mapa bidimensional de sensibilidad variando ligeramente dos parámetros iniciales (por ejemplo, anomalía media y velocidad radial inicial) dentro de rangos físicamente realistas. Dibujar curvas de nivel de la distancia mínima Tierra--Apophis y marcar zonas donde cambios diminutos en condiciones iniciales producen diferencias grandes en el resultado. Esto permite mostrar de manera visual qué tan ``delicada'' es la geometría del encuentro.

\subsection{Ventana de observación realista con clima y brillo del cielo}

Ir más allá de la geometría celeste e incorporar factores observacionales reales: fase lunar, altura sobre el horizonte, brillo del cielo urbano y probabilidad histórica de nubosidad para distintas ciudades de Colombia. Con ello, construir una ``ventana de observación efectiva'' por ciudad y proponer cuál sería el mejor sitio y hora para observación pública del evento.

\subsection{Reconstrucción inversa del encuentro}

Plantear el problema inverso: suponer que solo se conocen tres observaciones sintéticas de Apophis (ascensión recta y declinación en horas diferentes) y reconstruir, a partir de ellas, un estado inicial orbital aproximado. Comparar ese estado estimado con el estado real usado en la simulación y cuantificar el error en la predicción de la distancia mínima. La idea muestra cómo pequeñas campañas de observación pueden traducirse en predicciones dinámicas.

\subsection{Índice de riesgo didáctico para pasos cercanos futuros}

Diseñar un índice simple y transparente que combine tres variables: distancia mínima, velocidad relativa y tiempo dentro de la esfera de influencia terrestre. Aplicar ese índice al paso de 2029 y a uno o dos pasos futuros de Apophis para comparar escenarios. El objetivo no es reemplazar escalas oficiales, sino crear una métrica pedagógica que ayude a interpretar de forma física la ``intensidad'' de cada encuentro cercano.
